\documentclass[aspectratio=169,10pt]{beamer}
\usetheme{metropolis}
\usepackage{booktabs,amsmath,amssymb}
\title{Q14: Concepts, Features, and Memory}
\subtitle{Directions, superposition, and where knowledge lives}
\author{DLS 2026 --- Kukla}
\date{October 2026}
\begin{document}
\maketitle

\begin{frame}{The question}
\begin{block}{Student question (\#27)}
How are ``concepts'' maintained and generated? Do they live in features? How do models reason and keep
memory?
\end{block}
\end{frame}
\begin{frame}{Concepts = directions, not neurons}
\begin{itemize}
\item A concept $\approx$ a \textbf{direction} in activation space.
\item One neuron participates in \emph{many} concepts.
\item Linear representation hypothesis: features are directions; operations approx.\ linear.
\end{itemize}
\end{frame}
\begin{frame}{Superposition}
\begin{itemize}
\item Networks pack \textbf{more features than dimensions}.
\item Works when features are \emph{sparse} (few active at once).
\item Explains \textbf{polysemantic neurons}.
\item Sparse autoencoders disentangle them back into meaningful features.
\end{itemize}
\end{frame}
\begin{frame}{Where knowledge lives: MLP layers}
\begin{itemize}
\item Feed-forward sublayers act as \textbf{key--value memories}.
\item Causal tracing can localize --- and \emph{edit} --- a specific fact.
\item Attention \emph{moves} info; MLPs \emph{store} it.
\end{itemize}
\end{frame}
\begin{frame}{Two kinds of memory}
\begin{itemize}
\item \textbf{Weights}: long-term, learned in training, fixed at inference.
\item \textbf{Context}: short-term working memory, gone when the chat ends.
\end{itemize}
\alert{A model doesn't ``remember'' past chats unless re-supplied / retrieved (Q15).}
\end{frame}

\end{document}
